\begin{table}[!htbp]
\centering
\caption{Rótulos de bom desempenho e acurácia balanceada em amostra.}
\label{tab:res-rotulos}
\begin{singlespace}
\normalsize
\setlength{\tabcolsep}{5pt}
\renewcommand{\arraystretch}{1.15}
\begin{tabular}{lrrrr}
\hline\hline
\textbf{Método} & \textbf{\shortstack{Bons na\\coleção 55}} & \textbf{\shortstack{Bons na\\coleção 74}} & \textbf{\shortstack{BA controle 74\\(\%)}} & \textbf{\shortstack{BA topologia 74\\(\%)}} \\
\hline
IBS & 51/55 & 70/74 & 50,00 & 50,00 \\
MIP & 24/55 & 42/74 & 96,06 & 93,68 \\
CUTH & 7/55 & 24/74 & 93,83 & 92,83 \\
LBBD & 21/55 & 36/74 & 94,59 & 94,74 \\
ILS & 21/55 & 35/74 & 93,30 & 93,59 \\
RS & 6/55 & 23/74 & 94,88 & 97,83 \\
\hline\hline
\end{tabular}
\par\smallskip
{\small Nota: bom desempenho corresponde a \textit{gap} médio de até 10\%. BA é a acurácia balanceada. As métricas das duas últimas colunas usam as próprias instâncias do ajuste, sem avaliação externa.\par}
{\fonte{elaboração própria a partir dos dados experimentais.}}
\end{singlespace}
\end{table}
